\documentclass[11pt]{article}
\usepackage[margin=1.1in]{geometry}
\usepackage{amsmath,amssymb,amsthm,mathtools,booktabs,array,tabularx,longtable}
\usepackage[dvipsnames]{xcolor}
\usepackage[colorlinks=true,linkcolor=MidnightBlue,citecolor=MidnightBlue,urlcolor=MidnightBlue]{hyperref}
\usepackage{fancyhdr}
\pagestyle{fancy}\fancyhf{}\renewcommand{\headrulewidth}{0pt}
\fancyfoot[R]{\thepage}
\fancyfoot[L]{\scriptsize\color{gray}Generated by the Transport Lab from approved parts of its store; not a deposited paper. Papers of record: \href{https://zenodo.org/communities/transport-papers}{Zenodo, Transport Papers}}
\fancypagestyle{plain}{\fancyhf{}\fancyfoot[R]{\thepage}\fancyfoot[L]{\scriptsize\color{gray}Generated by the Transport Lab from approved parts of its store; not a deposited paper. Papers of record: \href{https://zenodo.org/communities/transport-papers}{Zenodo, Transport Papers}}\renewcommand{\headrulewidth}{0pt}}
\newtheorem{theorem}{Theorem}[section]
\newtheorem{proposition}[theorem]{Proposition}
\newtheorem{lemma}[theorem]{Lemma}
\newtheorem{corollary}[theorem]{Corollary}
\newtheorem{conjecture}[theorem]{Conjecture}
\newtheorem{refutation}[theorem]{Refutation}
\newtheorem{prediction}[theorem]{Prediction}
\newtheorem{principle}[theorem]{Principle}
\newtheorem{observation}[theorem]{Observation}
\newtheorem{axiom}[theorem]{Axiom}
\theoremstyle{definition}
\newtheorem{definition}[theorem]{Definition}
\newtheorem{fact}[theorem]{Fact}
\newtheorem{assembled}[theorem]{Assembly}
\newtheorem{postulate}[theorem]{Postulate}
\newtheorem{premise}[theorem]{Premise}
\newtheorem{hypothesis}[theorem]{Hypothesis}
\newtheorem{remark}[theorem]{Remark}
\newtheorem{example}[theorem]{Example}
\theoremstyle{plain}
\setlength{\marginparwidth}{0.8in}
\newcommand{\changedmark}[2]{\marginpar{\raggedright\scriptsize\color{#1}#2}}
\providecommand{\Aa}{}\renewcommand{\Aa}{\mathcal{A}}
\providecommand{\Ad}{}\renewcommand{\Ad}{\widetilde{\mathrm{Ad}}}
\providecommand{\Aev}{}\renewcommand{\Aev}{\mathcal{A}^{+}}
\providecommand{\Aod}{}\renewcommand{\Aod}{\mathcal{A}^{-}}
\providecommand{\B}{}\renewcommand{\B}{\mathrm{B}}
\providecommand{\C}{}\renewcommand{\C}{\mathbb{C}}
\providecommand{\CAR}{}\renewcommand{\CAR}{\mathrm{CAR}}
\providecommand{\CS}{}\renewcommand{\CS}{{\mathbb{S}^{+}}}
\providecommand{\CSm}{}\renewcommand{\CSm}{{\mathbb{S}^{-}}}
\providecommand{\Cl}{}\renewcommand{\Cl}{\mathrm{Cl}}
\providecommand{\Con}{}\renewcommand{\Con}{\mathcal{C}}
\providecommand{\Dl}{}\renewcommand{\Dl}{\Delta}
\providecommand{\EN}{}\renewcommand{\EN}{E_{N}}
\providecommand{\End}{}\renewcommand{\End}{\mathrm{End}}
\providecommand{\F}{}\renewcommand{\F}{\mathbb{F}}
\providecommand{\G}{}\renewcommand{\G}{\Gamma}
\providecommand{\HH}{}\renewcommand{\HH}{\mathbb{H}}
\providecommand{\IS}{}\renewcommand{\IS}{I\sigma_3}
\providecommand{\J}{}\renewcommand{\J}{\mathbf{J}}
\providecommand{\LH}{}\renewcommand{\LH}{L_{H}}
\providecommand{\Mat}{}\renewcommand{\Mat}{\mathrm{M}}
\providecommand{\Mo}{}\renewcommand{\Mo}{M^{\circ}}
\providecommand{\Mp}{}\renewcommand{\Mp}{M^{+}}
\providecommand{\Msun}{}\renewcommand{\Msun}{M_{\odot}}
\providecommand{\Mt}{}\renewcommand{\Mt}{\widetilde{M}}
\providecommand{\Ncone}{}\renewcommand{\Ncone}{\mathcal{N}}
\providecommand{\Neff}{}\renewcommand{\Neff}{N_{\mathrm{eff}}}
\providecommand{\PCS}{}\renewcommand{\PCS}{\mathbb{P}{\mathbb{S}^{+}}}
\providecommand{\PT}{}\renewcommand{\PT}{\mathbb{PT}}
\providecommand{\Pin}{}\renewcommand{\Pin}{\mathrm{Pin}}
\providecommand{\Q}{}\renewcommand{\Q}{\mathbb{Q}}
\providecommand{\R}{}\renewcommand{\R}{\mathbb{R}}
\providecommand{\Ric}{}\renewcommand{\Ric}{\mathrm{Ric}}
\providecommand{\SBH}{}\renewcommand{\SBH}{S_{\mathrm{BH}}}
\providecommand{\Sc}{}\renewcommand{\Sc}{\mathbb{S}}
\providecommand{\Sch}{}\renewcommand{\Sch}{\mathrm{P}}
\providecommand{\Sig}{}\renewcommand{\Sig}{\Sigma}
\providecommand{\Som}{}\renewcommand{\Som}{S_\omega}
\providecommand{\Sp}{}\renewcommand{\Sp}{\mathbb{S}}
\providecommand{\Sph}{}\renewcommand{\Sph}{\mathrm{S}}
\providecommand{\Spi}{}\renewcommand{\Spi}{S_\pi}
\providecommand{\Spin}{}\renewcommand{\Spin}{\mathrm{Spin}}
\providecommand{\Sth}{}\renewcommand{\Sth}{S_{\mathrm{th}}}
\providecommand{\T}{}\renewcommand{\T}{\mathbb{T}}
\providecommand{\TH}{}\renewcommand{\TH}{T_{\mathrm{H}}}
\providecommand{\Tr}{}\renewcommand{\Tr}{\operatorname{Tr}}
\providecommand{\UHF}{}\renewcommand{\UHF}{\mathrm{UHF}}
\providecommand{\Z}{}\renewcommand{\Z}{\mathbb{Z}}
\providecommand{\Zc}{}\renewcommand{\Zc}{\mathrm{Z}}
\providecommand{\cmg}{}\renewcommand{\cmg}{\mathrm{cm^{2}\,g^{-1}}}
\providecommand{\dS}{}\renewcommand{\dS}{\mathrm{dS}}
\providecommand{\dsc}{}\renewcommand{\dsc}{\mathrm{dsc}}
\providecommand{\dtPG}{}\renewcommand{\dtPG}{dt}
\providecommand{\dtS}{}\renewcommand{\dtS}{dt_{\mathrm{S}}}
\providecommand{\ellv}{}\renewcommand{\ellv}{\ell}
\providecommand{\er}{}\renewcommand{\er}{e_r}
\providecommand{\et}{}\renewcommand{\et}{e_t}
\providecommand{\ev}{}\renewcommand{\ev}{\mathrm{ev}}
\providecommand{\gl}{}\renewcommand{\gl}{\mathfrak{gl}}
\providecommand{\hatV}{}\renewcommand{\hatV}{\widehat{V}}
\providecommand{\hb}{}\renewcommand{\hb}{\bar{h}}
\providecommand{\hf}{}\renewcommand{\hf}{\hat{f}}
\providecommand{\id}{}\renewcommand{\id}{\mathrm{id}}
\providecommand{\kpc}{}\renewcommand{\kpc}{\mathrm{kpc}}
\providecommand{\lP}{}\renewcommand{\lP}{\ell_{P}}
\providecommand{\mD}{}\renewcommand{\mD}{\mathbf{m}_{D}}
\providecommand{\mb}{}\renewcommand{\mb}{\bar m}
\providecommand{\mnu}{}\renewcommand{\mnu}{\Sigma m_{\nu}}
\providecommand{\nInf}{}\renewcommand{\nInf}{n_\infty}
\providecommand{\nO}{}\renewcommand{\nO}{n_0}
\providecommand{\ninf}{}\renewcommand{\ninf}{n_{\infty}}
\providecommand{\nz}{}\renewcommand{\nz}{n_{0}}
\providecommand{\proj}{}\renewcommand{\proj}{\operatorname{proj}}
\providecommand{\pv}{}\renewcommand{\pv}{\pi}
\providecommand{\rad}{}\renewcommand{\rad}{\mathrm{rad}}
\providecommand{\sgn}{}\renewcommand{\sgn}{\operatorname{sgn}}
\providecommand{\shA}{}\renewcommand{\shA}{\mathcal{A}}
\providecommand{\shC}{}\renewcommand{\shC}{\mathcal{C}}
\providecommand{\shE}{}\renewcommand{\shE}{\mathcal{E}}
\providecommand{\shJ}{}\renewcommand{\shJ}{\mathcal{J}}
\providecommand{\shK}{}\renewcommand{\shK}{\mathcal{K}}
\providecommand{\shL}{}\renewcommand{\shL}{L}
\providecommand{\shW}{}\renewcommand{\shW}{\mathcal{W}}
\providecommand{\sll}{}\renewcommand{\sll}{\mathfrak{sl}}
\providecommand{\so}{}\renewcommand{\so}{\mathfrak{so}}
\providecommand{\su}{}\renewcommand{\su}{\mathfrak{su}}
\providecommand{\supp}{}\renewcommand{\supp}{\operatorname{supp}}
\providecommand{\tr}{}\renewcommand{\tr}{\operatorname{tr}}
\providecommand{\uu}{}\renewcommand{\uu}{\mathfrak{u}}
\providecommand{\br}{}\renewcommand{\br}[1]{\langle #1\rangle}
\providecommand{\bra}{}\renewcommand{\bra}[1]{\langle #1|}
\providecommand{\file}{}\renewcommand{\file}[1]{\texttt{#1}}
\providecommand{\hp}{}\renewcommand{\hp}[2]{( #1 \,|\, #2 )}
\providecommand{\ip}{}\renewcommand{\ip}[2]{\langle #1, #2\rangle}
\providecommand{\ket}{}\renewcommand{\ket}[1]{|#1\rangle}
\providecommand{\rel}{}\renewcommand{\rel}[2]{D\bigl(#1 \,\|\, #2\bigr)}
\providecommand{\rev}{}\renewcommand{\rev}[1]{\widetilde{#1}}
\providecommand{\slashed}{}\renewcommand{\slashed}[1]{{#1}\mkern-10mu/}
\title{\bfseries A tour of the Transport Papers}
\author{Leon Butler\\ \normalsize\href{mailto:lbutler.math@gmail.com}{\texttt{lbutler.math@gmail.com}}\\ \normalsize ORCID: \href{https://orcid.org/0009-0006-4449-5617}{0009-0006-4449-5617}}
\date{Generated by the Transport Lab, 2 November 2026}
\begin{document}
\maketitle
\begin{abstract}
\normalsize
A null vector, nonzero but with square zero, can never be divided by, even inside a larger algebra. With both space and time directions and at least three dimensions, no choice of partial inverses respects the symmetries, however the algebra is enlarged.

The fix is a larger space with one new positive and one new negative direction, $\Cl(p,q) \to \Cl(p{+}1, q{+}1)$. This becomes a recipe sorting every failure of division into three kinds, each with its own fix. Each step up the resulting tower adds exactly one fermionic mode.

At the centre of a non-rotating black hole the obstruction is the kind this repair supplies. Einstein's vacuum equations become one null tractor staying constant. The boundary modes carry a thermal state at the Hawking temperature, the only quasi-free state thermal for the boost clock.

On the matter side, the smallest massive twistor composite carries exactly the symmetry of the strong, weak and hypercharge forces. With an assumed hierarchical second correction, the neutrino masses add up to about 0.059 eV. Given stated symmetry assumptions, a mirror world copies the Standard Model.
\end{abstract}
\newpage
\noindent \textbf{Authorship and disclosure. }The founding idea of the series is the author's: that division by zero in a Clifford algebra should be read not as a prohibition but as a transport into a higher-dimensional algebra, and that the transported division need not be accompanied by an inverse operation. The direction of investigation and the critical review of the results are likewise the author's. The mathematical development and the drafting of the text were carried out in extended collaboration with Claude Fable 5 (Anthropic), an AI system.

\noindent\fbox{\parbox{\dimexpr\linewidth-2\fboxsep-2\fboxrule}{\textbf{This is not a deposited paper.} It was generated by the Transport Lab on 2 November 2026 and assembled from approved parts of its store. It follows results across several papers of the series; each is stated with its proof and credited to the paper that first proved it. The official versions of those papers are on Zenodo, in the Transport Papers community, \url{https://zenodo.org/communities/transport-papers}.}}
\newpage

\section*{Introduction}
This tour follows one idea from pure algebra all the way to particles, neutrinos and dark matter. It starts with a failure of division, because every later step is a response to that failure.

It then shows how enlarging the space repairs the failure, and applies the repair to geometry and gravity, then to black holes and heat, and last to matter. The assumptions that some results rest on are gathered first, under their own heading, so that what is proved stays separate from what is supposed.

\section{What we assume}

\begin{postulate}[P1: the single frozen point]\label{c:dark-matter.p1-single-frozen-point}
This assumption says the origin of the universe happened once, in one way, for everything. Every hidden sector feels the same gravity, and every sector that can have a handedness received the same one. It is assumed, not derived.
\end{postulate}

\begin{postulate}[P2: indifference, scoped to sector labels]\label{c:dark-matter.p2-indifference}
This assumption says the birth of the universe treats alike only things that are genuinely alike. The only hidden sectors that count as alike are pairs related by a reflection, so only those pairs are forced to match. Which member of the matching pair we live in is simply a matter of where we stand. What must match are the sector's fixed labels, such as its stored charge, not how hot it is: the assumption leaves the temperature of the mirror sector open, and it takes the number of modes the birth deposits to be finite.
\end{postulate}

\begin{postulate}[H: the second stage is hierarchical]\label{c:neutrino.h-second-stage-hierarchical}
This is an assumption, not a result: the lightest neutrino gets almost no mass, less than a hundredth of the next one. Other particle families show steep mass hierarchies of this kind, which makes the assumption plausible without proving it. With it, the framework's total neutrino mass becomes a sharp number rather than a lower limit.
\end{postulate}

\begin{postulate}[M: the framework realises the Majorana term, at a scale $M \gg v$]\label{c:neutrino.m-framework-realises-majorana-term-scale}
Neutrinos could carry a special kind of mass that only electrically neutral particles can have. Nothing forbids it, but nothing in the framework yet shows that it forms or how big it is. The paper assumes that it does, at a very high scale, and leaves the explanation for later work.
\end{postulate}

\begin{postulate}[Tameness]\label{c:sheaf.tameness}
This is the standing assumption on how a geometry is allowed to degenerate. The places where it degenerates must form a smooth closed wall, exactly one direction must lose its length at each point of the wall, and these directions must vary smoothly along it. In addition, the length must vanish at a linear rate as the wall is reached, not to higher order.
\end{postulate}

\section{Zero divisors and obstructions}

The series starts from a simple question: what happens when you try to divide by something that is not zero but whose square is zero?

In a space of mixed signature, $pq \neq 0$, like the spacetime of relativity, such vectors exist, and they form a cone much like a light cone.

The first result shows how firmly these vectors resist division.

A null $v \neq 0$ has no inverse on either side, in $\Cl(p,q)$ or in any larger algebra containing it.

\begin{lemma}[A null vector has no inverse]\label{c:transport.noinverse}
A null vector is a nonzero direction whose square is zero. Such an element can never be divided by: if it had an inverse on either side, multiplying through would force the vector itself to be zero. The argument uses only the fact that its square is zero, so it still applies in any larger algebra of the same kind that contains the original one.
\emph{First proved in} T1, Lem 2.1 (\href{https://doi.org/10.5281/zenodo.21743863}{doi:10.5281/zenodo.21743863}).
\end{lemma}
\noindent\emph{Rests on:} \ref{c:std-algebra.vector-squares}.

With exact inverses ruled out, one might settle for partial ones.

A partial inverse $\sigma(v)$ satisfies \[ v\,\sigma(v)\,v = v \] but in mixed signature with $n \geq 3$ no choice of one for each null $v$ is equivariant, even under $\Spin^0(p,q)$, in any extension.

\begin{theorem}[No-go at the level of elements]\label{c:transport.no-go-level-elements}
A light-like direction has no inverse, but it has many partial inverses. This theorem shows that there is no way to choose one for every light-like direction that respects the symmetries of the geometry, in any space with both space and time directions and at least three dimensions. Enlarging the algebra by any number of new directions does not help.
\emph{First proved in} T1, Thm 4.2 (\href{https://doi.org/10.5281/zenodo.21743863}{doi:10.5281/zenodo.21743863}).
\end{theorem}
\noindent\emph{Rests on:} \ref{c:transport.splitting-along-hyperbolic-plane}, \ref{c:primer.commutation-rule}, \ref{c:std-algebra.witt-decomposition}.

\section{Clifford algebra structure}

A Clifford algebra is built from a space of directions together with a rule for multiplying them. Before trying to repair these algebras, it helps to know what they really are.

The classical answer is a short and complete list.

Each real Clifford algebra is a matrix algebra over $\R$, $\C$ or $\HH$, or a sum of two, according to the residue \[ d \equiv p - q \pmod 8 \] and it is simple except when $p - q \equiv 1 \pmod 4$.

\begin{theorem}[Classification of the real Clifford algebras]\label{c:transport.classification-real-clifford-algebras}
Clifford algebras built from a space with p positive and q negative directions turn out to be familiar objects in disguise: square tables of numbers (matrices) whose entries are real numbers, complex numbers or quaternions, or a pair of such tables side by side. Which of these appears depends only on p minus q, in a pattern that repeats every eight steps. Usually the algebra is a single indivisible block; it splits into two identical pieces exactly when p minus q leaves remainder 1 on division by four.
\emph{First proved in} T1, Thm 2.5 (\href{https://doi.org/10.5281/zenodo.21743863}{doi:10.5281/zenodo.21743863}).
\end{theorem}
\noindent\emph{Rests on:} \ref{c:primer.classification-real-clifford-algebras}.

\section{The completion}

Since division cannot be rescued inside the algebra, the series moves outside it, into a larger space. The question for this section is how much larger that space must be, and what to do once one is there.

The first answer concerns the reciprocal of a single vector.

The larger space needs one new direction of each sign, \[ \Cl(p,q) \;\longrightarrow\; \Cl(p{+}1, q{+}1) \] and this hyperbolic plane is enough and unique up to isometry.

\begin{theorem}[Minimality]\label{c:transport.minimality}
To make the reciprocal work everywhere, the series represents each point as a light-like direction in a larger space. This result asks how much larger that space has to be. The added part must contain at least one direction of positive square and one of negative square, so adding only space-like or only time-like directions never works, while one such pair is enough and is essentially unique.
\emph{First proved in} T1, Thm 4.13 (\href{https://doi.org/10.5281/zenodo.21743863}{doi:10.5281/zenodo.21743863}).
\end{theorem}
\noindent\emph{Rests on:} \ref{c:std-algebra.witt-decomposition}, \ref{c:std-algebra.sylvester-inertia}.

A similar count applies when the geometry itself is degenerate.

A radical of dimension $r$ needs $r$ new partners, so $\Cl(p,q,r)$ embeds minimally in $\Cl(p{+}r, q{+}r)$.

\begin{theorem}[Minimal nondegenerate completion]\label{c:conformal.minimal-nondegenerate-completion}
A geometric algebra built on a degenerate metric has directions that pair to zero with everything, and division fails along them. This result says how to enlarge the space so that the metric becomes nondegenerate: each degenerate direction needs exactly one new partner direction, and the enlarged metric gains one positive and one negative direction for each degenerate one. The smallest such enlargement is unique up to a symmetry that leaves the original space untouched. Minimality here is among enlargements of the underlying space, not among all algebras.
\emph{First proved in} T2, Thm 2.3 (\href{https://doi.org/10.5281/zenodo.22714886}{doi:10.5281/zenodo.22714886}).
\end{theorem}
\noindent\emph{Rests on:} \ref{c:conformal.blade-basis-persists-under-degeneracy}, \ref{c:conformal.witt-extension-theorem}, \ref{c:std-algebra.functoriality}.

These pieces now fit together into one working procedure.

The theorems distil into a six-step procedure, and every step names the theorem that licenses it.

\begin{theorem}[The transport algorithm]\label{c:transport.transport-algorithm}
This is the series' recipe for what to do when division fails. Find where it fails, decide which of three kinds of failure it is, apply the matching fix (add a pair of dimensions, add one dimension, or do nothing), do the calculation in the enlarged setting, and bring the answer back, noting where it cannot come back. The recipe can be repeated, but its bookkeeping settles after at most one change.
\emph{First proved in} T1, Sec 8.1 (\href{https://doi.org/10.5281/zenodo.21743863}{doi:10.5281/zenodo.21743863}).
\end{theorem}
\noindent\emph{Rests on:} \ref{c:transport.classification-real-clifford-algebras}, \ref{c:transport.cohomological-form-no-go}, \ref{c:transport.descent-read-out-map}, \ref{c:transport.door-one-way}, \ref{c:transport.exhaustiveness-every-zero-divisor-annihilated-idempotent}, \ref{c:transport.existence-classification-type-c}, \ref{c:transport.existence-locus-type-n}, \ref{c:transport.fusion-occurs-most-once}, \ref{c:transport.minimality}, \ref{c:transport.no-fourth-rule}, \ref{c:transport.no-go-level-elements}, \ref{c:transport.no-transport-terminates}, \ref{c:transport.noinverse}, \ref{c:transport.one-generator-minimal-minimal-completions-classified}, \ref{c:transport.pinching-kernel}, \ref{c:transport.representation-independent-minimality-completion}, \ref{c:transport.rule-1-morita-trivial}, \ref{c:transport.rule-2-ideal-fusion-sharp-form}, \ref{c:transport.swap-rescale-normal-form}, \ref{c:transport.totality}, \ref{c:transport.type-p-self-repairing}.

\section{The tower and quantum capacity}

Repeating the basic repair builds a tower of ever larger spaces.

Each step adds a null pair $v_i, w_i$, and these obey exactly the rules of fermionic modes, $v_i w_j + w_j v_i = \delta_{ij}$.

The next result counts how much of this quantum structure each floor can hold.

The number of modes realisable by vectors is the Witt index, \[ \min(p,q) \] and along the tower it is $\min(p,q) + k$, one more per rung.

\begin{theorem}[Witt index is quantum capacity]\label{c:tower.witt-index-quantum-capacity}
Inside a Clifford algebra one can ask how many independent fermionic modes can be built directly from vectors. The answer is a classical number attached to the geometry, the Witt index: the largest number of independent zero-length directions that are all orthogonal to one another. For a space with p positive and q negative directions it is the smaller of p and q. Each step up the tower raises this number by exactly one, so each step adds precisely one mode of capacity and no more.
\emph{First proved in} T3, Thm 4.3 (\href{https://doi.org/10.5281/zenodo.22721556}{doi:10.5281/zenodo.22721556}).
\end{theorem}
\noindent\emph{Rests on:} \ref{c:tower.ladder-pairs-satisfy-car-vectors}, \ref{c:std-algebra.vector-squares}, \ref{c:std-algebra.witt-decomposition}.

\section{Degenerate metrics and walls}

Real geometries can degenerate in some places and not in others, so the repair has to be made locally, along a wall.

The standing tameness assumption fixes what is allowed: a smooth closed wall, exactly one direction losing its length at each point, and that length vanishing at a linear rate.

With that assumption in place, the repair can be carried right up to the wall.

On the physical side the completed form is nondegenerate of signature $(p{+}1, q{+}1)$ at every point, including the wall, so every fibre is \[ \Cl(2,4) \cong \Mat\_4(\HH) \] in the physical case.

\begin{theorem}[One-sided completion]\label{c:sheaf.one-sided-completion}
Assume a geometry degenerates along a wall in the tame way, with at least one positive direction on the physical side and with the dying direction's line consistently orientable along the wall. Then it can be repaired on that side by adding two extra dimensions and coupling them to the direction that is dying. This result shows that the coupling can be chosen so that the enlarged geometry is non-degenerate, with the same signature everywhere on the physical side, the wall itself included. The resulting family of algebras is then uniform: in the physical case every member is the algebra of four-by-four quaternion matrices.
\emph{First proved in} T4, Thm 4.4 (\href{https://doi.org/10.5281/zenodo.22721669}{doi:10.5281/zenodo.22721669}).
\end{theorem}
\noindent\emph{Rests on:} \ref{c:sheaf.determinant-identity-blow-up-expansion}, \ref{c:sheaf.local-freeness-constant-rank}, \ref{c:sheaf.wall-gram}, \ref{c:sheaf.tameness}.

\section{The singularity}

The centre of a black hole is the first physical place where the repair is put to work.

Approaching it, the gravitational field grows without bound, but once that growth is divided out the field settles to a limit that is degenerate.

The question is which kind of obstruction this degeneration turns out to be.

It is type N, a null direction without a partner, and type C cannot occur, because \[ p - q = -2 \equiv 6 \pmod 8 \] makes $\Cl(1,3) \cong \Mat_2(\HH)$ simple.

\begin{proposition}[Type N, and only type N]\label{c:singularity.type-n-only-type-n}
At the centre of a non-rotating black hole one direction of the geometry collapses, and this result identifies which kind of algebraic obstruction that is in the series' classification. It is the kind caused by a direction with no partner to pair with. The other substantive kind, which would need the algebra to split into independent pieces, cannot occur for the algebra of spacetime, and a third kind is present but harmless. So the repair that applies is the one that supplies the missing partner by enlarging the algebra.
\emph{First proved in} T5, Prop 3.1 (\href{https://doi.org/10.5281/zenodo.22721790}{doi:10.5281/zenodo.22721790}).
\end{proposition}
\noindent\emph{Rests on:} \ref{c:sheaf.localised-transport}, \ref{c:sheaf.one-sided-completion}, \ref{c:sheaf.one-sided-completion-singularity}, \ref{c:sheaf.schwarzschild-degeneration-boundary-tame}, \ref{c:singularity.degeneration-exactly}, \ref{c:tower.rigidity-clock-frozen}, \ref{c:transport.existence-classification-type-c}, \ref{c:transport.existence-locus-type-n}, \ref{c:transport.type-p-self-repairing}, \ref{c:primer.classification-real-clifford-algebras}, \ref{c:std-algebra.sta-matrix}.

\section{Einstein's equations in the completion}

To say anything about gravity itself, the series needs Einstein's equations in a form that lives naturally in the enlarged space.

The next result supplies such a form.

In Gauge Theory Gravity on flat spacetime, the vacuum field equations hold exactly when the null scale tractor $I_g = \nInf$ is parallel for the $\Cl(2,4)$ gauge field.

\begin{theorem}[Covariant form of the vacuum equations]\label{c:interior.covariant-form-vacuum-equations}
In the interior region, Einstein's equations for empty space can be recast as a single condition on the gravitational fields: a particular marker in the six-dimensional conformal extension of spacetime, attached to the physical scale, must stay constant as it is carried from point to point by a gauge field built from gravity. This result proves that the two formulations are exactly equivalent, and that the marker is a null (lightlike) object. It also notes that merely having some constant marker of this kind is a weaker condition that does not by itself give the field equations.
\emph{First proved in} T6, Thm 4.5 (\href{https://doi.org/10.5281/zenodo.22721990}{doi:10.5281/zenodo.22721990}).
\end{theorem}
\noindent\emph{Rests on:} \ref{c:interior.imported-equivalence-einstein-equations}, \ref{c:interior.imported-tractor-connection}, \ref{c:interior.metric-parallelism-trace-schouten-condition}.

\section{Descent, parity and states}

After working upstairs one has to come back down, and any return journey must lose something. This section says what is lost and what the loss means.

The first result identifies the way down with a familiar operation.

Descent is the average over fermion parity, \[ E(x) \;=\; \tfrac12\bigl(x + \alpha(x)\bigr) \] whose kernel is exactly the adjoined ladder sector.

\begin{theorem}[The descent is the parity average]\label{c:singularity.descent-parity-average}
The descent, which keeps a mode's occupation data and discards its coherences, is the same thing as averaging over fermion parity: take an observable and its parity-flipped version and average the two. This makes it the only consistent projection onto the parity-even observables, and it leaves every parity-even state unchanged.
\emph{First proved in} T5, Thm 5.7 (\href{https://doi.org/10.5281/zenodo.22721790}{doi:10.5281/zenodo.22721790}).
\end{theorem}
\noindent\emph{Rests on:} \ref{c:singularity.descent-what-it-erases}.

That leaves the question of where the missing partner of a mixed state comes from.

In the state's own representation the partner is the commutant, and for a Gibbs state the cyclic vector is the thermofield double.

\begin{theorem}[Purity is an output: the partner is the commutant]\label{c:singularity.purity-output-partner-commutant}
Starting from a single quantum mode in a mixed state, there is a standard construction that builds a larger system in which that state is one half of a pure entangled state. This result shows that the partner needed for the pure state is not an extra ingredient: it emerges automatically as the mirror copy of the mode's own algebra, and the construction is unique up to changes made on the partner side alone. For a thermal state the resulting pure state is the familiar thermofield double.
\emph{First proved in} T5, Thm 7.27 (\href{https://doi.org/10.5281/zenodo.22721790}{doi:10.5281/zenodo.22721790}).
\end{theorem}
\noindent\emph{Rests on:} \ref{c:singularity.modular-theory-matrix-factor-proved-outright}.

\section{Temperature and Hawking radiation}

The series next asks whether the modes it finds at the black-hole boundary can be hot, in the sense that Hawking described.

The thermal case of that construction now meets the black hole directly.

At inverse temperature \[ \beta = 1/\TH = 2\pi/\kappa \] the boundary modes have a state thermal for the boost clock, and within the quasi-free class it is unique, its two-point data forced.

\begin{theorem}[The I-state]\label{c:hawking.i-state}
The boundary modes in this black-hole setting can be placed in a thermal state at the Hawking temperature, each mode filled according to the Fermi-Dirac rule. This result shows that this state exists and is thermal in the precise algebraic sense, both on every finite set of modes and on the whole system. It also shows that among the simple Gaussian (quasi-free) states it is the only one that is thermal for the geometry's boost time evolution, which turns the phase of each mode at its own frequency, so its occupation numbers are derived rather than chosen.
\emph{First proved in} T7, Thm 3.6 (\href{https://doi.org/10.5281/zenodo.22722121}{doi:10.5281/zenodo.22722121}).
\end{theorem}
\noindent\emph{Rests on:} \ref{c:hawking.gibbs-states-are-kms-conversely-finite}, \ref{c:hawking.imported-kms-passes-inductive-limits}, \ref{c:hawking.imported-quasi-free-states-car-algebra}.

\section{Dark energy and the frozen point}

The cosmological constant enters through the frozen point, a fixed direction in an enlarged conformal space whose length the constant sets. One of the earlier results used here rests on an assumed principle rather than a proof.

From heat the tour turns to gravity on the largest scale.

Where $\Lambda \neq 0$, the frozen point supplies the Stelle-West vector, so gravity's action makes sense where the metric degenerates.

\begin{theorem}[The frozen point makes the Stelle-West vector]\label{c:dark-energy.frozen-point-makes-stelle-west-vector}
The frozen point is a fixed direction in the larger conformal space whose length encodes the cosmological constant. On the region where Einstein's equations hold with a nonzero cosmological constant, this result shows that the frozen point, together with one canonical geometric object, produces exactly the extra vector that a known reformulation of gravity, due to MacDowell and Mansouri and to Stelle and West, requires. Gravity's action can then be written as the square of the curvature as seen from the point at infinity, with no inverse metric anywhere, so it is defined wherever the connection stays regular. At the black hole's wall, where the connection is singular, its value diverges but its field equations still reach the wall; and the curvature-squared form itself exists only while the cosmological constant is nonzero.
\emph{First proved in} T10, Thm 3.1 (\href{https://doi.org/10.5281/zenodo.22722496}{doi:10.5281/zenodo.22722496}).
\end{theorem}
\noindent\emph{Rests on:} \ref{c:dark-energy.imported-frozen-point-geometry}, \ref{c:dark-energy.imported-macdowell-mansouri-theorem-stelle-west}, \ref{c:dark-energy.imported-walls-sphere-theorem}, \ref{c:dark-energy.stelle-west-vector-frozen-point}.

\section{The origin of mass}

The tour now turns from gravity to matter, using twistors, in which a single twistor describes a massless particle travelling at light speed.

The first question about matter is where its mass comes from.

A single conformal spinor carries the null momentum \[ p\_{AA'} = \pi\_{A}\bar\pi\_{A'} \] and is massless, so mass first appears for a system of two.

\begin{proposition}[Mass as symplectic area: the two-spinor formula]\label{c:mass.mass-symplectic-area-two-twistor-formula}
A single spinor of the conformal geometry describes a massless particle travelling at light speed. Two of them together can describe a massive particle, and its mass depends only on how far apart their two light-like directions point. When they line up the composite stays massless; the more they disagree, the heavier it is.
\emph{First proved in} T11, Prop 4.1 (\href{https://doi.org/10.5281/zenodo.22722565}{doi:10.5281/zenodo.22722565}).
\end{proposition}

\section{The Standard Model}

Next the tour asks which forces such a massive composite would feel.

Looking inside the smallest massive composite answers that question.

Two conformal spinors of three pairs each have automorphisms $(U(2) \times U(3))/U(1)$, with Lie algebra \[ \su(2) \oplus \su(3) \oplus \uu(1) \] the Standard Model's exactly.

\begin{theorem}[The gauge algebra of the minimal massive composite]\label{c:standard-model.gauge-algebra-minimal-massive-composite}
A composite of basic spinor states of the conformal algebra has a two-level internal structure: which spinor state, and which of its three internal pairs. The rotations that respect both levels form a symmetry algebra. For the smallest composite that can have mass, two spinor states of three pairs each, that algebra is exactly the one behind the strong, weak and hypercharge forces; for three spinor states it is too large.
\emph{First proved in} T12, Thm 4.1 (\href{https://doi.org/10.5281/zenodo.22722669}{doi:10.5281/zenodo.22722669}).
\end{theorem}
\noindent\emph{Rests on:} \ref{c:mass.mass-symplectic-area-two-twistor-formula}.

\section{Families of quarks and leptons}

The next question is why matter might repeat itself in families.

Here Hypothesis (S), an assumption about how the slots of a composite behave, is needed, and several of the framework's particle-physics results hold only if it does. In particular, everything said about charged leptons depends on it.

Allowing a composite more than two twistors leads to a simple count of choices.

The faces, three of them under the three-spinor bound, carry identical content, because the permutations $S_{3}$ of the slots carry each face's copy onto every other.

\begin{proposition}[Three identical families, mixing by shared slots]\label{c:families.three-identical-families-mixing-shared-slots}
If the largest allowed composite is made of three conformal spinors and a massive particle uses two of them, there are three ways to choose the pair. Each choice carries an exact copy of the same set of particles, because relabelling the spinors maps one copy onto another; the copies include the lepton doublet if the slots behave like the pairs (Hypothesis (S)). Any two choices share exactly one spinor, which hints at why quark families mix the way they do, though that hint and the matching of choices to real families are not proved.
\emph{First proved in} T13, Prop 3.2 (\href{https://doi.org/10.5281/zenodo.22722845}{doi:10.5281/zenodo.22722845}).
\end{proposition}
\noindent\emph{Rests on:} \ref{c:families.imported-one-derived-generation-hypothesis-s}, \ref{c:families.imported-face-valued-bracket-quantum-numbers}, \ref{c:standard-model.hyp-s-slot-statistics}.

\section{Neutrino masses}

Neutrinos are where the framework comes closest to numbers that experiments can check.

The account assumes that neutrinos carry a Majorana mass at a very high scale, which nothing in the framework yet shows forming. The family picture it builds on is itself still a conjecture.

Treating the families even-handedly has a sharp consequence for neutrino masses.

The Dirac mass matrix is democratic, $\mD = d\,\J$, with the single eigenvalue $m_{D} = 3d$, so the seesaw gives the normal ordering.

\begin{proposition}[The rank-one seesaw: ordering, sum, and the double-beta values]\label{c:neutrino.rank-one-seesaw-ordering-sum-double}
Assuming the framework supplies the heavy neutrino's Majorana mass at a scale far above the electroweak one (Premise (M)), and because the framework's neutrino couplings treat the three families identically at first approximation (a structural-level reading), only one light neutrino gets a mass, through the seesaw. Provided the corrections that break that symmetry are small, that fixes the mass ordering as normal; if the second stage of those corrections is also hierarchical (Premise (H)), the total mass sits at the smallest value oscillation experiments allow, about 0.059 eV, and otherwise that value is only a lower bound. Under the same premises neutrinoless double-beta decay occurs, at rates too small for the next generation of experiments to see: a conditional absence.
\emph{First proved in} T15, Prop 4.4 (\href{https://doi.org/10.5281/zenodo.22723070}{doi:10.5281/zenodo.22723070}).
\end{proposition}
\noindent\emph{Rests on:} \ref{c:neutrino.imported-indifference-faces-democratic-alignment-rank}, \ref{c:neutrino.invariant-symmetric-matrices}, \ref{c:neutrino.seesaw-what-gauge-invariance-proves-what}, \ref{c:neutrino.h-second-stage-hierarchical}, \ref{c:neutrino.m-framework-realises-majorana-term-scale}.

\section{Dark matter and the dark towers}

The account of matter ends with matter we cannot see. The answer rests on postulates, not proofs, about how the origin of the universe treats its hidden sectors.

Swapping the roles of time and space gives a second world to set beside our own.

Every input of the Standard Model chain is invariant under the reversal \[ (p,q) \mapsto (q,p) \] which carries $(1,3)$ to the mirror's $(3,1)$, so the mirror has the same gauge algebra and the same generation.

\begin{theorem}[The mirror tower is a copy of this one]\label{c:dark-matter.mirror-tower-copy-one}
Swapping which directions count as time and which as space turns our world into a mirror world, and this result shows that every ingredient used to derive our particle physics survives the swap. So the mirror world has its own copy of the Standard Model as the series derives it (under Hypothesis (S)), with mirror photons and mirror neutrinos. If the two worlds share one cutoff scale (a common ceiling), then under the framework's premises of a single frozen point (P1) and indifference (P2), which fix a common coupling there and an equal matter excess, it also has the same particle masses and the same matter excess, so its matter equals our ordinary matter to within half a percent, the gap coming from the two worlds' nuclear histories. A later note (T14c) shows that the same indifference principle, applied as widely, would also make the mirror world as hot as ours, which observations rule out, so the cold mirror world the series relies on is an unstated exception to that principle.
\emph{First proved in} T14, Thm 4.3 (\href{https://doi.org/10.5281/zenodo.22722953}{doi:10.5281/zenodo.22722953}).
\end{theorem}
\noindent\emph{Rests on:} \ref{c:dark-matter.imported-inputs-standard-model-chain}, \ref{c:dark-matter.type-residue-what-each-tower-can}, \ref{c:dark-matter.p1-single-frozen-point}, \ref{c:dark-matter.p2-indifference}.

\section{One relation, read twice}

The tour closes by stepping back to the single rule behind everything.

That rule, $\{e_{a}, e_{b}\} = 2g_{ab}\mathbf{1}$, can be read as geometry or as the rule obeyed by fermions, though only at the level of structure.

A last result counts the quantum structure in two ways.

The real and complex capacities are \[ r\_{\R}(p,q) = \min(p,q), \qquad r\_{\C}(p,q) = \lfloor d/2 \rfloor \] and for spacetime $(1,3)$ they are $1$ and $2$.

\begin{proposition}[The two capacities, and three consequences]\label{c:one-relation.two-capacities-three-consequences}
Reading spacetime geometry as a system of fermion modes, there are two ways to count the modes: using only real directions, or allowing complex ones. The two counts usually differ, and for ordinary spacetime they give one and two. Each step up the tower of extended spaces adds exactly one mode to both counts, and the complex count of the extended six-dimensional space is three.
\emph{First proved in} T16, Prop 3.2 (\href{https://doi.org/10.5281/zenodo.22723212}{doi:10.5281/zenodo.22723212}).
\end{proposition}
\noindent\emph{Rests on:} \ref{c:one-relation.imported-witts-extension-theorem}, \ref{c:one-relation.two-presentations}, \ref{c:tower.witt-index-quantum-capacity}, \ref{c:std-algebra.vector-squares}, \ref{c:std-algebra.witt-decomposition}.

\section{Results used from elsewhere}
The proofs above use the following results of the series, stated briefly; each is proved in the paper named.

\begin{fact}[The blade basis persists under degeneracy]\label{c:conformal.blade-basis-persists-under-degeneracy}
For any symmetric bilinear form on an $m$-dimensional real space, degenerate or not, the $2^m$ blades of any basis form a linear basis of $\Cl(W,B)$.
\emph{Source:} T2, Fact 2.1 (\href{https://doi.org/10.5281/zenodo.22714886}{doi:10.5281/zenodo.22714886}).
\end{fact}

\begin{fact}[The Witt extension theorem]\label{c:conformal.witt-extension-theorem}
Over a field of characteristic not two, an isometry between arbitrary subspaces of two isometric nondegenerate quadratic spaces extends to an isometry of the whole spaces.
\emph{Source:} T. Y. Lam, Introduction to Quadratic Forms over Fields, Graduate Studies in Mathematics 67, American Mathematical Society (2005)., Ch.~I, Thm.~4.7.
\end{fact}

\begin{fact}[Imported: the frozen point and its geometry]\label{c:dark-energy.imported-frozen-point-geometry}
On the invertible region of a gravitational field on flat spacetime the Einstein condition makes the scale tractor parallel for the tractor gauge field, with $Q(I) = \Lambda/3$, stabiliser $SO(1,4)$ or $SO(2,3)$, and $P_{ab} = (\Lambda/6)g_{ab}$.
\emph{Source:} T10, Fact 2.1 (\href{https://doi.org/10.5281/zenodo.22722496}{doi:10.5281/zenodo.22722496}).
\end{fact}

\begin{fact}[Imported: the MacDowell-Mansouri theorem, Stelle-West form]\label{c:dark-energy.imported-macdowell-mansouri-theorem-stelle-west}
The MacDowell-Mansouri functional of a de Sitter gauge field $A(a) = \Omega(a) - \hb^{-1}(a)\wedge V/\ell^2$ on flat spacetime, with a Stelle-West vector of fixed norm, splits into Gauss-Bonnet, Einstein-Cartan and cosmological terms, and its field equations are the Einstein-Cartan equations with cosmological constant $\Lambda$.
\emph{Source:} S. W. MacDowell and F. Mansouri, Unified geometric theory of gravity and supergravity, Phys. Rev. Lett. 38 (1977), 739--742.; K. S. Stelle and P. C. West, Spontaneously broken de Sitter symmetry and the gravitational holonomy group, Phys. Rev. D 21 (1980), 1466--1488.; D. K. Wise, MacDowell-Mansouri gravity and Cartan geometry, Class. Quantum Grav. 27 (2010), 155010..
\end{fact}

\begin{fact}[Imported: walls, and the sphere theorem]\label{c:dark-energy.imported-walls-sphere-theorem}
At a tame corank-one wall the Clifford sheaf continues, the metric has a simple pole with $\mathcal{L}^{\otimes 2} \cong N^*\Sigma$, the completion adjoins a Witt partner of the radical covector, and, under the Indifference principle, the transverse sphere is carried by boundary-mode bilinears.
\emph{Source:} T10, Fact 2.4 (\href{https://doi.org/10.5281/zenodo.22722496}{doi:10.5281/zenodo.22722496}).
\end{fact}

\begin{fact}[The Stelle-West vector of the frozen point]\label{c:dark-energy.stelle-west-vector-frozen-point}
In the frozen point's scale, $V = X - (3/\Lambda)I$ is orthogonal to $I$ with $Q(V) = -3/\Lambda$ and $\nabla V = e$, and the tractor gauge field is $\Omega(a) + \hb^{-1}(a)\wedge V/Q(V)$.
\emph{Source:} T10, Lem 2.5 (\href{https://doi.org/10.5281/zenodo.22722496}{doi:10.5281/zenodo.22722496}).
\end{fact}

\begin{fact}[Imported: the inputs of the Standard-Model chain]\label{c:dark-matter.imported-inputs-standard-model-chain}
The Standard-Model chain of the series uses exactly a stated list of inputs, from the conformal completion and its conformal spinor ideal to the mass formula, pair-counting hypercharge, Hypothesis (S) and signature-free imports, a list carried as a hypothesis.
\emph{Source:} T14, Fact 4.2 (\href{https://doi.org/10.5281/zenodo.22722953}{doi:10.5281/zenodo.22722953}).
\end{fact}

\begin{fact}[Type by residue, and what each tower can carry intrinsically]\label{c:dark-matter.type-residue-what-each-tower-can}
A tower's spinor type (real, complex or quaternionic, with one or two modules) is a function of its residue alone, and it fixes the intrinsic symmetry group, $\Z_{2}$, $U(1)$ or $Sp(1)$, that each tower's ground algebra carries.
\emph{Source:} T14, Thm 3.4 (\href{https://doi.org/10.5281/zenodo.22722953}{doi:10.5281/zenodo.22722953}).
\end{fact}

\begin{fact}[Imported: the face-valued bracket and its quantum numbers]\label{c:families.imported-face-valued-bracket-quantum-numbers}
The mass bracket transforms by $\det(U)$ under the internal $U(2)$ and carries the quantum numbers $(1,1)$ with $Y = +1$ of $e^{c}$, not those of a weak doublet.
\emph{Source:} T13, Fact 2.1 (\href{https://doi.org/10.5281/zenodo.22722845}{doi:10.5281/zenodo.22722845}).
\end{fact}

\begin{fact}[Imported: one derived generation, and Hypothesis (S)]\label{c:families.imported-one-derived-generation-hypothesis-s}
The quark doublet and the singlet hypercharges follow unconditionally, while the lepton doublet, and, given the one-handed content and $Y_{Q} = 1/6$, the anomaly-fixed uniqueness of the generation's hypercharges, hold only under Hypothesis (S).
\emph{Source:} T13, Fact 2.3 (\href{https://doi.org/10.5281/zenodo.22722845}{doi:10.5281/zenodo.22722845}).
\end{fact}

\begin{fact}[Gibbs states are KMS, and conversely, in finite dimensions]\label{c:hawking.gibbs-states-are-kms-conversely-finite}
On a full matrix algebra with dynamics generated by a diagonal Hamiltonian, the Gibbs state at inverse temperature $\beta$ is KMS, and it is the unique KMS state at that $\beta$, degenerate spectrum included.
\emph{Source:} T7, Lem 3.1 (\href{https://doi.org/10.5281/zenodo.22722121}{doi:10.5281/zenodo.22722121}).
\end{fact}

\begin{fact}[Imported: KMS passes to inductive limits]\label{c:hawking.imported-kms-passes-inductive-limits}
A state that is KMS at $\beta$ on each subalgebra of an increasing chain of dynamics-invariant subalgebras is KMS at $\beta$ on the closure of their union.
\emph{Source:} O. Bratteli and D. W. Robinson, Operator Algebras and Quantum Statistical Mechanics 2, 2nd ed., Springer, Berlin (1997)., Prop.~5.3.25.
\end{fact}

\begin{fact}[Imported: quasi-free states of the CAR algebra]\label{c:hawking.imported-quasi-free-states-car-algebra}
Each symbol $0 \le A \le 1$ determines a unique gauge-invariant quasi-free state of the CAR algebra; quasi-free states are fixed by their symbol and pairing, and gauge invariance is the vanishing of the pairing.
\emph{Source:} H. Araki, On quasifree states of CAR and Bogoliubov automorphisms, Publ. RIMS Kyoto Univ. 6 (1971), 385--442.; O. Bratteli and D. W. Robinson, Operator Algebras and Quantum Statistical Mechanics 2, 2nd ed., Springer, Berlin (1997)., \S 5.2.
\end{fact}

\begin{fact}[Imported: equivalence with the Einstein equations]\label{c:interior.imported-equivalence-einstein-equations}
Wherever $\hb$ is invertible, the GTG field equations are equivalent to Einstein's equations for the induced metric, and in vacuum to its Ricci-flatness.
\emph{Source:} A. Lasenby, C. Doran and S. Gull, Gravity, gauge theories and geometric algebra, Phil. Trans. R. Soc. Lond. A 356 (1998), 487--582., \S 4 (\S 4.3, covariant forms of the field equations), App.~C; C. Doran and A. Lasenby, Geometric Algebra for Physicists, Cambridge University Press (2003)., Ch.~13-14.
\end{fact}

\begin{fact}[Imported: the tractor connection]\label{c:interior.imported-tractor-connection}
For a gravitational field on flat spacetime, the tractor bundle of its conformal class (the $\Cl(2,4)$ vectors $\sigma\nInf + \mu + \rho\nO$) carries a canonical metric-preserving connection, given in a chosen scale by an explicit formula involving the Schouten tensor.
\emph{Source:} T. N. Bailey, M. G. Eastwood and A. R. Gover, Thomas's structure bundle for conformal, projective and related structures, Rocky Mountain J. Math. 24 (1994), 1191--1217.; S. N. Curry and A. R. Gover, An introduction to conformal geometry and tractor calculus, with a view towards analysis on manifolds, in Asymptotic Analysis in General Relativity, London Mathematical Society Lecture Note Series 443, Cambridge University Press (2018), 86--170., \S 5.
\end{fact}

\begin{fact}[Metric parallelism, and the trace of the Schouten condition]\label{c:interior.metric-parallelism-trace-schouten-condition}
The tractor gauge field preserves the tractor metric, and a Schouten map proportional to the identity forces the factor to be $R/24$ and the field to be Einstein, so $\Sch = 0$ exactly when $\Ric = 0$.
\emph{Source:} T6, Lem 4.2 (\href{https://doi.org/10.5281/zenodo.22721990}{doi:10.5281/zenodo.22721990}).
\end{fact}

\begin{fact}[Imported: indifference, the faces, the democratic alignment, and the rank-one texture]\label{c:neutrino.imported-indifference-faces-democratic-alignment-rank}
Under the indifference premise, the conjectural face picture of the families and the Gram-square reading, the leading Yukawa form is democratic, proportional to the all-ones matrix, giving one heavy and two light families at leading order.
\emph{Source:} T15, Fact 2.3 (\href{https://doi.org/10.5281/zenodo.22723070}{doi:10.5281/zenodo.22723070}).
\end{fact}

\begin{fact}[$S_{3}$-invariant symmetric matrices]\label{c:neutrino.invariant-symmetric-matrices}
The permutation-invariant real symmetric $3 \times 3$ matrices are exactly $\alpha\mathbf{1} + \beta\J$, with eigenvalue $\alpha + 3\beta$ on $u'$ and $\alpha$ on $u'^{\perp}$.
\emph{Source:} T15, Lem 4.3 (\href{https://doi.org/10.5281/zenodo.22723070}{doi:10.5281/zenodo.22723070}).
\end{fact}

\begin{fact}[The seesaw: what gauge invariance proves, and what is assumed]\label{c:neutrino.seesaw-what-gauge-invariance-proves-what}
The cosmological term cannot give the neutrino mass, the singlet's Majorana term is the one mass gauge invariance allows without a condensate, and under (M) the leading-order one-generation seesaw gives $m_{\nu} = m_{D}^{2}/M_{R}$.
\emph{Source:} T15, Prop 4.1 (\href{https://doi.org/10.5281/zenodo.22723070}{doi:10.5281/zenodo.22723070}).
\end{fact}

\begin{fact}[Imported: Witt's extension theorem]\label{c:one-relation.imported-witts-extension-theorem}
Over a field of characteristic not two, every isometry between subspaces of a nondegenerate quadratic space extends to the whole space, so the isometry group is transitive on subspaces of each isometry type.
\emph{Source:} E. Witt, Theorie der quadratischen Formen in beliebigen K"orpern, J. Reine Angew. Math. 176 (1937) 31--44.; T. Y. Lam, Introduction to Quadratic Forms over Fields, Graduate Studies in Mathematics 67, American Mathematical Society (2005)., Ch.~I, \S4.
\end{fact}

\begin{fact}[The two presentations]\label{c:one-relation.two-presentations}
Read on an orthogonal basis the Clifford relation is a metric, and read on a Witt basis it is the canonical anticommutation relations, so the geometric and quantum readings are two charts on one algebra.
\emph{Source:} T16, Prop 3.1 (\href{https://doi.org/10.5281/zenodo.22723212}{doi:10.5281/zenodo.22723212}).
\end{fact}

\begin{fact}[Classification of the real Clifford algebras]\label{c:primer.classification-real-clifford-algebras}
$\Cl(p,q)$ is a matrix algebra over $\R$, $\C$ or $\HH$, or a sum of two such, according to $p - q \bmod 8$, and it is simple except when $p - q \equiv 1 \pmod 4$.
\emph{Source:} T0, Thm 9.3 (\href{https://doi.org/10.5281/zenodo.22714501}{doi:10.5281/zenodo.22714501}).
\end{fact}

\begin{fact}[The commutation rule]\label{c:primer.commutation-rule}
If $x$ lies in the subalgebra generated by vectors orthogonal to $v$, then $v x = \hat{x} v$ and $x v = v \hat{x}$.
\emph{Source:} T0, Lem 9.1 (\href{https://doi.org/10.5281/zenodo.22714501}{doi:10.5281/zenodo.22714501}).
\end{fact}

\begin{fact}[The determinant identity, and the blow-up expansion]\label{c:sheaf.determinant-identity-blow-up-expansion}
Off the wall within the collar, $\det h = \det G \cdot \bigl(-(1 + 2\mu\nu\beta^2 s)\bigr)$ exactly, and $s = 1/\lambda + O(1)$ diverges with the sign of $\lambda$ as the wall is approached.
\emph{Source:} T4, Lem 4.3 (\href{https://doi.org/10.5281/zenodo.22721669}{doi:10.5281/zenodo.22721669}).
\end{fact}

\begin{fact}[Local freeness at constant rank]\label{c:sheaf.local-freeness-constant-rank}
For any cometric, degenerate or not, the Clifford sheaf is locally free of constant rank $2^n$ with smooth structure functions, so its rank never detects degeneration.
\emph{Source:} T4, Prop 2.8 (\href{https://doi.org/10.5281/zenodo.22721669}{doi:10.5281/zenodo.22721669}).
\end{fact}

\begin{fact}[Localised transport]\label{c:sheaf.localised-transport}
The degenerate Clifford sheaf embeds canonically in the completed sheaf, and on the wall the radical covector acquires a partner with which it forms a ladder pair that exists exactly over the wall.
\emph{Source:} T4, Cor 4.5 (\href{https://doi.org/10.5281/zenodo.22721669}{doi:10.5281/zenodo.22721669}).
\end{fact}

\begin{fact}[One-sided completion at the singularity]\label{c:sheaf.one-sided-completion-singularity}
At the blown-up Schwarzschild singularity, coupling an adjoined hyperbolic plane to the dying direction with $\mu\nu < 0$ and bump inside the horizon gives an ambient form of constant signature $(2,4)$ everywhere, boundary included, with no shadow wall and a ladder pair existing exactly over the boundary.
\emph{Source:} T4, Thm 8.12 (\href{https://doi.org/10.5281/zenodo.22721669}{doi:10.5281/zenodo.22721669}).
\end{fact}

\begin{fact}[The Schwarzschild degeneration is boundary-tame in $\shW$]\label{c:sheaf.schwarzschild-degeneration-boundary-tame}
In the weighted module at the blown-up Schwarzschild singularity, the derived cometric of the Painlev\'e-Gullstrand gauge field extends smoothly to the boundary and degenerates there with corank exactly one and to exactly first order, with radical $k = \rho + m\,\dtPG$ annihilating the infalling congruence.
\emph{Source:} T4, Thm 8.7 (\href{https://doi.org/10.5281/zenodo.22721669}{doi:10.5281/zenodo.22721669}).
\end{fact}

\begin{fact}[The wall Gram]\label{c:sheaf.wall-gram}
On the wall the ambient form splits orthogonally into the surviving directions and a nondegenerate three-dimensional coupled block, giving signature $(p{+}1, q{+}1)$ when $\mu\nu > 0$ and $(p, q{+}2)$ when $\mu\nu < 0$.
\emph{Source:} T4, Lem 4.2 (\href{https://doi.org/10.5281/zenodo.22721669}{doi:10.5281/zenodo.22721669}).
\end{fact}

\begin{fact}[The degeneration, exactly]\label{c:singularity.degeneration-exactly}
As $r \to 0$ the rescaled gauge map and metric converge exactly to rank-one limits whose common kernel $\et^{\perp}$ contains the folded direction $\er$, null for the limit form.
\emph{Source:} T5, Prop 2.6 (\href{https://doi.org/10.5281/zenodo.22721790}{doi:10.5281/zenodo.22721790}).
\end{fact}

\begin{fact}[The descent, and what it erases]\label{c:singularity.descent-what-it-erases}
The descent $E(x) = PxP + (1-P)x(1-P)$, with $P = v^* v$, keeps a mode's diagonal occupation data and erases its coherences.
\emph{Source:} T5, Prop 5.5 (\href{https://doi.org/10.5281/zenodo.22721790}{doi:10.5281/zenodo.22721790}).
\end{fact}

\begin{fact}[Modular theory of a matrix factor, proved outright]\label{c:singularity.modular-theory-matrix-factor-proved-outright}
For a full matrix algebra with a cyclic separating vector, the modular operator, modular conjugation and modular flow are computed explicitly, and the state is KMS for that flow.
\emph{Source:} T5, Lem 7.1 (\href{https://doi.org/10.5281/zenodo.22721790}{doi:10.5281/zenodo.22721790}).
\end{fact}

\begin{fact}[Hypothesis (S): slot statistics, derived from the slot-module hypothesis]\label{c:standard-model.hyp-s-slot-statistics}
Under the slot-module hypothesis, the composite's slots are fermionic modes obeying exclusion and forming one Clifford system with its pairs, its state space is the five-mode Fock space $\Lambda^{\bullet}(\C^{2} \oplus W)$ with $U(2) \times U(3)$ acting as $\Lambda^{\bullet}(A \oplus B)$, chirality is the parity of the mode number, and slots of one tower parity whose grading is the chirality are odd.
\emph{Source:} T12, Hypothesis (S) (\href{https://doi.org/10.5281/zenodo.22722669}{doi:10.5281/zenodo.22722669}).
\end{fact}

\begin{fact}[Isometries extend to algebra homomorphisms]\label{c:std-algebra.functoriality}
An isometry $f$ extends uniquely to an algebra homomorphism $\Cl(f)$, injective when $f$ is; $\mathrm{O}(V, Q)$ acts by automorphisms.
\emph{Source:} the algebra core of the Transport Lab, proved there with an exact check; standard, see P. Lounesto, Clifford Algebras and Spinors, 2nd ed., London Mathematical Society Lecture Note Series 286, Cambridge University Press (2001), Ch. 14 (definitions of the Clifford algebra).
\end{fact}

\begin{fact}[The spacetime algebra as a matrix algebra, and the spacetime split]\label{c:std-algebra.sta-matrix}
$\Cl(1,3) \cong \Mat_2(\HH)$, the Dirac matrices represent $\C \otimes \Cl(1,3) \cong \Mat_4(\C)$, and $\sigma_k = \gamma_k\gamma_0$ give $\Cl^{+}(1,3) \cong \Cl(3,0)$ with $\sigma_1\sigma_2\sigma_3 = I$.
\emph{Source:} the algebra core of the Transport Lab, proved there with an exact check; standard, see C. Doran and A. Lasenby, Geometric Algebra for Physicists, Cambridge University Press (2003), Ch. 5 (relativity and spacetime).
\end{fact}

\begin{fact}[Orthogonal bases and the law of inertia]\label{c:std-algebra.sylvester-inertia}
Every real quadratic space has an orthogonal basis squaring to $1, -1, 0$, and the counts $(p, q, r)$ are invariants that classify it up to isometry.
\emph{Source:} T. Y. Lam, Introduction to Quadratic Forms over Fields, AMS (2005)., the diagonalisation of forms and Sylvester's law of inertia (real case).
\end{fact}

\begin{fact}[The product of two vectors]\label{c:std-algebra.vector-squares}
$v^{2} = Q(v)$ and $uv + vu = 2B(u, v)$, so $uv = B(u, v) + \tfrac12(uv - vu)$.
\emph{Source:} the algebra core of the Transport Lab, proved there with an exact check; standard, see C. Doran and A. Lasenby, Geometric Algebra for Physicists, Cambridge University Press (2003), Ch. 4 (foundations of geometric algebra).
\end{fact}

\begin{fact}[Hyperbolic pairs and the Witt decomposition]\label{c:std-algebra.witt-decomposition}
A null vector in a nondegenerate plane has a unique hyperbolic partner, and $V = H_1 \perp \dots \perp H_{\min(p,q)} \perp A \perp \operatorname{rad} V$ with $A$ definite; maximal totally null subspaces have dimension $\min(p,q) + r$.
\emph{Source:} the algebra core of the Transport Lab, proved there with an exact check; standard, see T. Y. Lam, Introduction to Quadratic Forms over Fields, Graduate Studies in Mathematics 67, American Mathematical Society (2005), Ch. I (hyperbolic planes, the Witt decomposition, Witt's theorems).
\end{fact}

\begin{fact}[The ladder pairs satisfy the CAR, as vectors]\label{c:tower.ladder-pairs-satisfy-car-vectors}
In the $k$-th tower algebra the $k$ adjoined null pairs satisfy the canonical anticommutation relations of $k$ fermionic modes exactly, as vectors, and each new rung adds one mode without altering the existing relations.
\emph{Source:} T3, Prop 4.1 (\href{https://doi.org/10.5281/zenodo.22721556}{doi:10.5281/zenodo.22721556}).
\end{fact}

\begin{fact}[Rigidity; the clock is frozen]\label{c:tower.rigidity-clock-frozen}
Every floor of the tower shares the ground floor's charge, division-ring type, centre, simplicity class, Morita class and obstruction taxonomy, so the physical tower $\Cl(1{+}k, 3{+}k) \cong \Mat_{2^{k+1}}(\HH)$ is frozen at clock position $6$.
\emph{Source:} T3, Thm 3.6 (\href{https://doi.org/10.5281/zenodo.22721556}{doi:10.5281/zenodo.22721556}).
\end{fact}

\begin{fact}[Cohomological form of the no-go]\label{c:transport.cohomological-form-no-go}
The obstruction to an invariant pseudo-inverse of a null vector is the tautological class $\theta \in \mathrm{Hom}(N_v, U)$, which vanishes exactly when $U = 0$ and persists unchanged through every ascent and every larger symmetry group containing $N_v$.
\emph{Source:} T1, Thm 4.8 (\href{https://doi.org/10.5281/zenodo.21743863}{doi:10.5281/zenodo.21743863}).
\end{fact}

\begin{fact}[Descent: the read-out map]\label{c:transport.descent-read-out-map}
The read-out $\dsc[s\,\nO + x + t\,\nInf] = x/s$ inverts the lift off the ideal points, returns $x/Q(x)$ after the transported reciprocal, and is undefined exactly where $Q(x) = 0$.
\emph{Source:} T1, Lem 4.23 (\href{https://doi.org/10.5281/zenodo.21743863}{doi:10.5281/zenodo.21743863}).
\end{fact}

\begin{fact}[The door is one-way]\label{c:transport.door-one-way}
There is no unital algebra homomorphism from either minimal completion of a nondegenerate $\Cl(p,q)$ back to $\Cl(p,q)$ (the Rule 2 one existing only when $p - q \equiv 1 \pmod 4$), so every descent is non-multiplicative.
\emph{Source:} T1, Thm 7.6 (\href{https://doi.org/10.5281/zenodo.21743863}{doi:10.5281/zenodo.21743863}).
\end{fact}

\begin{fact}[Exhaustiveness: every zero divisor is annihilated by an idempotent]\label{c:transport.exhaustiveness-every-zero-divisor-annihilated-idempotent}
Every nonzero zero divisor of a nondegenerate real Clifford algebra is annihilated on each side by a nontrivial idempotent, so the N, P, C taxonomy of obstructions is exhaustive.
\emph{Source:} T1, Thm 3.9 (\href{https://doi.org/10.5281/zenodo.21743863}{doi:10.5281/zenodo.21743863}).
\end{fact}

\begin{fact}[Existence and classification of type C]\label{c:transport.existence-classification-type-c}
$\Cl(p,q)$, $n \geq 1$, has a nontrivial central idempotent exactly when $p - q \equiv 1 \pmod 4$, and then the only ones are $\varepsilon_\pm = \tfrac12(1 \pm \omega)$, which split the algebra into two simple ideals and are exchanged by the grade involution but by no inner automorphism.
\emph{Source:} T1, Thm 3.6 (\href{https://doi.org/10.5281/zenodo.21743863}{doi:10.5281/zenodo.21743863}).
\end{fact}

\begin{fact}[Existence locus of type N]\label{c:transport.existence-locus-type-n}
$\Cl(p,q)$ contains null vectors exactly when $pq \neq 0$; they then form a cone of dimension $n - 1$, and each is a two-sided zero divisor with no left or right inverse in any unital algebra containing $\Cl(p,q)$.
\emph{Source:} T1, Prop 3.3 (\href{https://doi.org/10.5281/zenodo.21743863}{doi:10.5281/zenodo.21743863}).
\end{fact}

\begin{fact}[Fusion occurs at most once]\label{c:transport.fusion-occurs-most-once}
Along any transport history Rule 1 fixes the mod-8 residue and Rule 2 leaves the type-C residues for good, so fusion happens at most once and the Wedderburn type is then fixed.
\emph{Source:} T1, Prop 7.4 (\href{https://doi.org/10.5281/zenodo.21743863}{doi:10.5281/zenodo.21743863}).
\end{fact}

\begin{fact}[No fourth rule]\label{c:transport.no-fourth-rule}
Every zero divisor of a nondegenerate real Clifford algebra is of type N, C or P, handled respectively by a hyperbolic-plane ascent, a one-generator ascent, or no ascent, so there is no fourth rule.
\emph{Source:} T1, Cor 6.7 (\href{https://doi.org/10.5281/zenodo.21743863}{doi:10.5281/zenodo.21743863}).
\end{fact}

\begin{fact}[No transport terminates]\label{c:transport.no-transport-terminates}
Every algebra reached by at least one transport step has zero divisors, of type N after Rule 1 and of type P after Rule 2, so transport never reaches a division algebra.
\emph{Source:} T1, Prop 7.8 (\href{https://doi.org/10.5281/zenodo.21743863}{doi:10.5281/zenodo.21743863}).
\end{fact}

\begin{fact}[One generator is minimal; the minimal completions classified]\label{c:transport.one-generator-minimal-minimal-completions-classified}
For $p - q \equiv 1 \pmod 4$, the minimal simple Clifford completions of $\Cl(p,q)$ are exactly the two one-generator ascents $\Cl(p{+}1,q)$ and $\Cl(p,q{+}1)$, isomorphic as algebras, while a null generator only thickens the algebra.
\emph{Source:} T1, Thm 5.7 (\href{https://doi.org/10.5281/zenodo.21743863}{doi:10.5281/zenodo.21743863}).
\end{fact}

\begin{fact}[The pinching, and its kernel]\label{c:transport.pinching-kernel}
The pinching by the descent projector is a conditional expectation onto its commutant whose kernel is exactly the adjoined ladder sector, and it returns transported elements unchanged.
\emph{Source:} T1, Prop 7.5 (\href{https://doi.org/10.5281/zenodo.21743863}{doi:10.5281/zenodo.21743863}).
\end{fact}

\begin{fact}[Representation-independent minimality of the completion]\label{c:transport.representation-independent-minimality-completion}
Any nondegenerate quadratic space whose null cone carries a normalised conformal chart of $(\mathbb{V}, Q)$ has $p'' \geq p + 1$ and $q'' \geq q + 1$, so the minimality of the completion does not depend on how $\Cl(p,q)$ is embedded.
\emph{Source:} T1, Thm 4.19 (\href{https://doi.org/10.5281/zenodo.21743863}{doi:10.5281/zenodo.21743863}).
\end{fact}

\begin{fact}[Rule 1 is Morita-trivial]\label{c:transport.rule-1-morita-trivial}
Rule 1 gives $\Cl(p+1, q+1) \cong \Mat_2(\Cl(p,q))$ with the adjoined ladder pair as off-diagonal matrix units, so it preserves the Morita class, the Wedderburn type and the clock residue $p - q \pmod 8$.
\emph{Source:} T1, Thm 7.1 (\href{https://doi.org/10.5281/zenodo.21743863}{doi:10.5281/zenodo.21743863}).
\end{fact}

\begin{fact}[Rule 2: ideal fusion, sharp form]\label{c:transport.rule-2-ideal-fusion-sharp-form}
For $p - q \equiv 1 \pmod 4$, adjoining one generator $\nu$ with $Q(\nu) = \pm 1$ gives the simple algebra $\Mat_2(\Cl^0(p,q))$ with centre $\R$, the same for either sign.
\emph{Source:} T1, Thm 5.5 (\href{https://doi.org/10.5281/zenodo.21743863}{doi:10.5281/zenodo.21743863}).
\end{fact}

\begin{fact}[Splitting along a hyperbolic plane]\label{c:transport.splitting-along-hyperbolic-plane}
Along a hyperbolic plane spanned by $(v, w_0)$, every element of $\Cl(H_0 \perp U')$ expands uniquely as $\alpha + v\beta + w_0\gamma + v w_0\delta$ with coefficients in $\Cl(U')$, and left multiplication by $v$ is injective on $\Cl(U')$.
\emph{Source:} T1, Lem 2.4 (\href{https://doi.org/10.5281/zenodo.21743863}{doi:10.5281/zenodo.21743863}).
\end{fact}

\begin{fact}[Swap-and-rescale normal form]\label{c:transport.swap-rescale-normal-form}
In homogeneous coordinates $(s, x, t)$ with $Q(x) = 2st$, the transported reciprocal is the linear involution $(s, x, t) \mapsto (2t, x, \frac12 s)$, and the ideal points $s = 0$ fall into three strata.
\emph{Source:} T1, Prop 4.21 (\href{https://doi.org/10.5281/zenodo.21743863}{doi:10.5281/zenodo.21743863}).
\end{fact}

\begin{fact}[Totality]\label{c:transport.totality}
On the quadric of the completion, $\iota[N] = [e_+ N e_+]$ is a total continuous involution that equals the reciprocal $x \mapsto x/Q(x)$ on invertible vectors and sends each $x_0$ with $Q(x_0) = 0$ to $[x_0 + \frac12 \nInf]$, and it is the only such continuous extension.
\emph{Source:} T1, Thm 4.10 (\href{https://doi.org/10.5281/zenodo.21743863}{doi:10.5281/zenodo.21743863}).
\end{fact}

\begin{fact}[Type P is self-repairing]\label{c:transport.type-p-self-repairing}
In any unital ring the identity assignment $\sigma(e) = e$ on idempotents satisfies both von Neumann identities and is equivariant under every automorphism of every extension preserving the ring, so type P obstructions are repaired at signature increment $(0,0)$.
\emph{Source:} T1, Prop 6.1 (\href{https://doi.org/10.5281/zenodo.21743863}{doi:10.5281/zenodo.21743863}).
\end{fact}

\begin{prediction}[Neutrino mass sum at the oscillation floor]\label{c:pred.neutrino-mass-sum-floor}
Under Premise (H) the neutrino mass sum lies at the normal-ordering floor, $\mnu \approx 0.059$ eV.
\emph{Kill condition:} Below: a physical-prior cosmological bound confirmed below the floor, or confirmed inverted ordering. Above: a sum determined meaningfully above the floor (near 0.075 eV), or m1 measured above a few meV.
\end{prediction}

\begin{prediction}[No neutrinoless double-beta signal at next-generation sensitivity]\label{c:pred.no-neutrinoless-double-beta}
Under Premise (H), $m_{\beta\beta}$ is a few meV, below next-generation double-beta sensitivity.
\emph{Kill condition:} A confirmed signal above the predicted window.
\end{prediction}

\begin{prediction}[The no-echo prediction, registered]\label{c:scrambling.no-echo-prediction-registered}
Under Postulate (C1), which puts the quantum structure at the stretched singularity, post-merger ringdown is the unmodified Kerr quasinormal spectrum, with no echo train at the delay $8M\ln(M/m_P)$.
\emph{Kill condition:} One statistically confirmed echo train at the Planck-structure delay.
\end{prediction}

\section*{Conclusion}
The thread of the tour is one move, made again and again. Division fails along null directions, so the space is enlarged by a pair of new directions, one of each sign. The work is done in the larger setting and the answer is brought back, noting where it cannot come back.

At the centre of a non-rotating black hole that move supplies the missing partner. Its boundary modes then carry a thermal state at the Hawking temperature whose occupation numbers are derived rather than chosen. On the matter side, the same algebra yields the symmetry of the strong, weak and hypercharge forces.

Not all of this is proved from nothing, and the assumptions are part of the picture. The repair along a wall assumes the degeneration is tame. The family results assume Hypothesis (S) about slot statistics.

The neutrino results assume that the framework realises a heavy Majorana mass, which is left for later work. The sharp total neutrino mass further assumes that the lightest neutrino gets almost no mass.

The mirror world assumes a single shared origin for every hidden sector. It also assumes that only sectors related by a reflection are forced to match. That the matter excess is imprinted equally in the two mirror sectors is proved for the mirror copy, not assumed.

Under these assumptions the framework makes claims that can fail. It predicts that the three neutrino masses add up to about 0.059 eV, with the normal ordering. It expects no neutrinoless double-beta signal at next-generation sensitivity, so a clear signal would count against it. It predicts no echoes after black holes merge, and a single confirmed echo would refute it.

\end{document}